\documentclass[12pt,a4paper]{article}
\usepackage[utf8]{inputenc}
\usepackage[T1]{fontenc}
\usepackage[italian]{babel}
\usepackage{geometry}
\usepackage{amsmath}
\usepackage{amssymb}
\usepackage{listings}
\usepackage{xcolor}
\usepackage{graphicx}
\usepackage{hyperref}
\usepackage{booktabs}
\usepackage{caption}
\usepackage{float}
\usepackage{parskip}
\usepackage{array}
\usepackage{tabularx}
\usepackage{tcolorbox}
\tcbuselibrary{listings,breakable,skins}

\geometry{
  top=2.5cm, bottom=2.5cm,
  left=2.5cm, right=2.5cm
}

\definecolor{codebg}{RGB}{248,248,248}
\definecolor{keywordcolor}{RGB}{0,0,180}
\definecolor{commentcolor}{RGB}{0,128,0}
\definecolor{stringcolor}{RGB}{163,21,21}
\definecolor{codeframe}{RGB}{180,180,180}

\lstdefinestyle{python}{
  language=Python,
  basicstyle=\ttfamily\footnotesize,
  keywordstyle=\color{keywordcolor}\bfseries,
  commentstyle=\color{commentcolor}\itshape,
  stringstyle=\color{stringcolor},
  numberstyle=\tiny\color{gray},
  numbers=left,
  stepnumber=1,
  breaklines=true,
  breakatwhitespace=false,
  showstringspaces=false,
  tabsize=4,
  xleftmargin=1.5em,
  keepspaces=true,
  columns=flexible,
  frame=none,
  aboveskip=0pt,
  belowskip=0pt,
}
\lstset{style=python}

\hypersetup{
  colorlinks=true,
  linkcolor=blue!60!black,
  urlcolor=blue!60!black,
}

\newtcblisting{pylisting}[1][]{
  listing only,
  enhanced,
  colback=codebg,
  colframe=codeframe,
  coltitle=black,
  fonttitle=\small\itshape,
  attach boxed title to top left={yshift=-2mm, xshift=4mm},
  boxed title style={colback=codebg, colframe=codeframe, boxrule=0.5pt},
  boxrule=0.5pt,
  arc=2pt,
  left=4pt, right=4pt, top=8pt, bottom=4pt,
  listing options={style=python},
  before upper={\vspace{-2pt}},
  #1
}

\newcommand{\HRule}{\rule{\linewidth}{0.5mm}}

\begin{document}
\raggedbottom

% ============================================================
% FRONTESPIZIO
% ============================================================
\begin{titlepage}
  \centering
  \vspace*{1.5cm}

  {\Large\textbf{Universit\`a degli Studi di Torino}\\[0.3em]
  Corso di Laurea Magistrale in Informatica}

  \vspace{1cm}
  \HRule\\[0.4em]
  {\LARGE\textbf{Tecnologie del Linguaggio Naturale}}\\[0.3em]
  {\large Parte 3 -- Prof.\ Luigi Di Caro}\\[0.2em]
  \HRule

  \vspace{1.2cm}
  {\huge\textbf{Lab 5}}\\[0.5em]
  {\LARGE Hybrid Search RAG}\\[0.4em]
  {\large Retrieval-Augmented Generation con ricerca ibrida semantica e lessicale}

  \vspace{2cm}

  {\large\textbf{Gruppo di lavoro:}}\\[0.8em]
  \begin{tabular}{ll}
    \toprule
    \textbf{Studente}    & \textbf{Matricola} \\
    \midrule
    Alessandro Olivero   & 915069  \\
    Samuele Iudice       & 1235245 \\
    Andrea Franco        & 1226739 \\
    \bottomrule
  \end{tabular}

  \vspace{1.5cm}
  \begin{tabular}{ll}
    \textbf{Docente:}          & Prof.\ Luigi Di Caro \\[0.3em]
    \textbf{Anno Accademico:}  & 2025/2026 \\
  \end{tabular}

  \vfill
\end{titlepage}

\tableofcontents
\newpage

% ============================================================
\section{Introduzione}
% ============================================================

Questo laboratorio affronta il problema della \textbf{Retrieval-Augmented Generation (RAG)}, un paradigma che potenzia i modelli di linguaggio combinando un meccanismo di \textit{retrieval} su una knowledge base esterna con la capacit\`a generativa di un Large Language Model (LLM). L'obiettivo \`e confrontare due strategie di retrieval su un corpus scientifico NLP:

\begin{itemize}
  \item \textbf{RAG base}: recupero puramente semantico tramite embeddings SciBERT e indice FAISS.
  \item \textbf{RAG ibrido}: combinazione del retrieval semantico con BM25 lessicale, fusi tramite Reciprocal Rank Fusion (RRF).
\end{itemize}

L'esperimento principale varia la dimensione del corpus da 1\,000 a 44\,949 documenti, misurando l'impatto sulle prestazioni con quattro metriche ispirate al framework RAGAS: Precision@5, Context Relevance, Answer Relevance e Faithfulness. Si analizza inoltre la \textbf{robustezza} del sistema a parafrasature delle query.

Il dataset usato \`e \texttt{MaartenGr/arxiv\_nlp}, una collezione di paper scientifici NLP da arXiv, disponibile su Hugging Face. Il modello di linguaggio \`e \textbf{Phi-3.5-mini-instruct} di Microsoft, eseguito interamente in locale su CPU.

% ============================================================
\section{Background Teorico}
% ============================================================

\subsection{Retrieval-Augmented Generation (RAG)}

Un sistema RAG si articola in tre componenti:
\begin{enumerate}
  \item \textbf{Retriever}: dato una query, recupera i documenti pi\`u rilevanti da una knowledge base.
  \item \textbf{Context assembly}: i documenti recuperati vengono concatenati come contesto.
  \item \textbf{Generator}: un LLM produce una risposta condizionata dal contesto e dalla query.
\end{enumerate}

Il vantaggio principale rispetto a un LLM stand-alone \`e la capacit\`a di accedere a informazioni aggiornate e specifiche di dominio senza ri-addestrare il modello, riducendo al contempo le allucinazioni.

\subsection{Retrieval semantico con FAISS e SciBERT}

Il retrieval semantico si basa su \textbf{dense embeddings}: ogni documento e ogni query vengono proiettati in uno spazio vettoriale denso, e il retrieval \`e una ricerca dei vicini pi\`u prossimi (Nearest Neighbor Search).

Si usa \textbf{SciBERT} (\texttt{allenai/scibert\_scivocab\_uncased}) come modello di embedding: un BERT pre-addestrato su 1,14 milioni di paper scientifici da Semantic Scholar. Produce vettori di dimensione 768, normalizzati a norma unitaria per l'uso con similarit\`a coseno.

\textit{Nota sulla scelta del modello}: SciBERT non \`e un modello \texttt{sentence-transformers} nativo, cio\`e non \`e stato fine-tuned con un obiettivo di similarit\`a tra frasi/documenti; caricandolo con \texttt{SentenceTransformer(...)} si ottiene un mean pooling di default (con relativo warning a runtime), che funziona ma pu\`o soffrire del problema noto di \textit{anisotropia} degli embedding BERT non specializzati per il retrieval. La scelta resta difendibile per il vocabolario specializzato nel dominio scientifico, ma un'alternativa pi\`u indicata sarebbe \texttt{allenai/specter} (o \texttt{specter2}), una variante di SciBERT gi\`a addestrata esplicitamente per la similarit\`a tra paper scientifici (tramite citazioni come segnale di supervisione).

L'indice di ricerca \`e costruito con \textbf{FAISS} (Facebook AI Similarity Search), che implementa \texttt{IndexFlatIP} (inner product su vettori normalizzati, equivalente alla similarit\`a coseno) e consente ricerca efficiente su decine di migliaia di vettori.

\subsection{Retrieval lessicale con BM25}

\textbf{BM25} (Best Match 25) \`e un modello di recupero probabilistico basato sulla frequenza dei termini. Per ogni query $q$ e documento $d$, il punteggio \`e:

\[
  \text{BM25}(q, d) = \sum_{t \in q} \text{IDF}(t) \cdot
  \frac{f(t,d) \cdot (k_1+1)}{f(t,d) + k_1 \left(1 - b + b \cdot \frac{|d|}{\text{avgdl}}\right)}
\]

dove $f(t,d)$ \`e la frequenza del termine $t$ nel documento $d$, $|d|$ \`e la lunghezza del documento, $\text{avgdl}$ \`e la lunghezza media del corpus, e $k_1 = 1.5$, $b = 0.75$ sono iperparametri standard.

A differenza del retrieval semantico, BM25 \`e efficace per query con termini tecnici specifici (es.\ acronimi, nomi propri) che gli embeddings tendono a generalizzare.

\subsection{Reciprocal Rank Fusion (RRF)}

La \textbf{Reciprocal Rank Fusion} \`e una tecnica di fusione dei risultati di pi\`u retriever. Dati i ranking prodotti da $n$ retriever, il punteggio RRF di un documento $d$ \`e:

\[
  \text{RRF}(d) = \sum_{r=1}^{n} \frac{1}{k + \text{rank}_r(d)}
\]

con $k=60$ (costante di smorzamento). I documenti vengono ri-ordinati per punteggio RRF decrescente. Questo approccio \`e robusto perch\'e non richiede normalizzazione degli score eterogenei e valorizza i documenti consistentemente rilevanti per pi\`u retriever.

\subsection{Phi-3.5-mini-instruct}

Il modello generativo \`e \textbf{Phi-3.5-mini-instruct} di Microsoft: un LLM da 3.8B parametri ottimizzato per il ragionamento e l'instruction-following. Viene eseguito in modalit\`a \textit{float32} su CPU, con generazione limitata a 300 token. Il template di chat segue il formato Phi-3 con ruoli \texttt{system} e \texttt{user}.

\subsection{Metriche di valutazione ispirate a RAGAS}

Si implementano quattro metriche per valutare il sistema RAG, ispirate al framework RAGAS (Retrieval-Augmented Generation Assessment System):

\begin{itemize}
  \item \textbf{Precision@5}: percentuale di documenti rilevanti (per gold standard) tra i top-5 recuperati.
  \[
    P@5 = \frac{|\{d \in R_5 : d \text{ rilevante}\}|}{5}
  \]

  \item \textbf{Context Relevance (CR)}: similarit\`a coseno media tra l'embedding SciBERT della query e gli snippet dei documenti recuperati. Misura quanto il contesto \`e semanticamente pertinente alla domanda.
  \[
    \text{CR} = \frac{1}{|R|} \sum_{d \in R} \cos(\text{enc}(q),\, \text{enc}(d[:200]))
  \]

  \item \textbf{Answer Relevance (AR)}: similarit\`a coseno tra l'embedding della query e l'embedding della risposta generata dall'LLM. Misura se la risposta \`e centrata sulla domanda.
  \[
    \text{AR} = \cos(\text{enc}(q),\, \text{enc}(a))
  \]

  \item \textbf{Faithfulness proxy (FF)}: percentuale di parole contenutistiche della risposta (escluse le stopwords) che compaiono nel testo dei documenti recuperati. Misura in modo approssimato quanto la risposta \`e fondata sul contesto.
  \[
    \text{FF} = \frac{|\{w \in a_{\text{content}} : w \in C\}|}{|a_{\text{content}}|}
  \]
  dove $C$ \`e il testo concatenato dei documenti di contesto.

  \item \textbf{Robustezza}: per un sottoinsieme di query, si calcolano le parafrasature e si misura l'overlap@5 tra i risultati della query originale e quelli delle parafrasature:
  \[
    \text{Rob} = \frac{1}{|P|} \sum_{p \in P} \frac{|R_5(q) \cap R_5(p)|}{5}
  \]
  Un valore alto indica che il retrieval \`e stabile a variazioni di formulazione della query.
\end{itemize}

% ============================================================
\section{Architettura del Sistema}
% ============================================================

Il sistema si struttura in tre livelli:

\begin{enumerate}
  \item \textbf{Livello dati}: caricamento del dataset arXiv-NLP, campionamento di $N$ documenti, costruzione del testo come concatenazione di titolo e abstract.
  \item \textbf{Livello retrieval}: indice FAISS (semantico) + indice BM25 (lessicale), con fusione RRF per il retrieval ibrido.
  \item \textbf{Livello generazione}: Phi-3.5-mini-instruct riceve la query e i top-5 documenti come contesto e genera una risposta in linguaggio naturale con citazioni.
\end{enumerate}

\begin{pylisting}[title={Funzione di retrieval ibrido con RRF}]
def retrieve_hybrid(query, faiss_index, bm25, documents, metadata, k=5):
    # retrieval semantico (top-20) e lessicale (top-20)
    sem = retrieve_semantic(query, faiss_index, documents, metadata, k=20)
    bm  = retrieve_bm25(query, bm25, documents, metadata, k=20)
    # fusione RRF e restituzione dei top-k
    return reciprocal_rank_fusion([sem, bm])[:k]

def reciprocal_rank_fusion(results_list, k=60):
    rrf_scores, doc_map = {}, {}
    for results in results_list:
        for item in results:
            ix = item['idx']
            rrf_scores[ix] = rrf_scores.get(ix, 0) + 1 / (k + item['rank'])
            doc_map[ix]    = item
    sorted_docs = sorted(rrf_scores.items(),
                         key=lambda x: x[1], reverse=True)
    return [
        {'rank': i+1, 'rrf_score': score, 'idx': ix,
         'document': doc_map[ix]['document'],
         'metadata': doc_map[ix]['metadata']}
        for i, (ix, score) in enumerate(sorted_docs)
    ]
\end{pylisting}

Il retrieval ibrido parte da $k'=20$ risultati per ciascun retriever prima della fusione, cos\`i da avere un pool di candidati sufficientemente ampio per il re-ranking.

% ============================================================
\section{Dataset e Configurazione}
% ============================================================

\subsection{Dataset}

Il dataset \texttt{MaartenGr/arxiv\_nlp} contiene \textbf{44\,949 paper scientifici NLP} da arXiv, con i campi \texttt{title}, \texttt{abstract} e \texttt{year}. Ogni documento \`e rappresentato come la concatenazione di titolo e abstract.

\begin{pylisting}[title={Caricamento e costruzione del corpus}]
dataset = load_dataset("MaartenGr/arxiv_nlp")
df_full = pd.DataFrame(dataset['train'])
# Documenti totali disponibili: 44,949

def build_corpus(df, n):
    df_s = df.sample(n=n, random_state=42).reset_index(drop=True)
    docs, meta = [], []
    for _, row in df_s.iterrows():
        title    = str(row['title'])
        abstract = str(row['abstract'])
        docs.append(f"Title: {title}\n\nAbstract: {abstract}")
        meta.append({'title': title, 'year': str(row['year'])})
    return docs, meta
\end{pylisting}

\subsection{Configurazione dell'esperimento}

L'esperimento varia la dimensione del corpus su cinque livelli:
\[
  N_{\text{DOCS}} \in \{1\,000,\; 5\,000,\; 15\,000,\; 25\,000,\; 44\,949\}
\]

Per ogni configurazione si ricostruisce il corpus (campionamento casuale con seed fisso), si riaddestramo gli indici FAISS e BM25, e si rieseguono le valutazioni. Il gold standard \`e definito da cinque query NLP con parole chiave rilevanti:

\begin{table}[H]
  \centering
  \caption{Query di valutazione e parole chiave del gold standard}
  \begin{tabular}{p{7cm} p{7cm}}
    \toprule
    \textbf{Query} & \textbf{Parole chiave rilevanti} \\
    \midrule
    What is BERT and how does self-attention work? & bert, attention, transformer, language model, pre-training \\
    How does word2vec represent word meaning?       & word2vec, skip-gram, word embedding, word representation \\
    What are the main challenges in machine translation? & machine translation, neural machine translation, nmt \\
    How does topic modeling work?                  & topic model, lda, bertopic, topic, clustering \\
    What is named entity recognition?              & named entity, ner, entity recognition, sequence labeling \\
    \bottomrule
  \end{tabular}
\end{table}

\subsection{Modelli utilizzati}

\begin{table}[H]
  \centering
  \caption{Modelli e librerie del sistema}
  \begin{tabular}{lll}
    \toprule
    \textbf{Componente}    & \textbf{Modello/Libreria}                      & \textbf{Ruolo} \\
    \midrule
    Embedding              & SciBERT (\texttt{allenai/scibert\_scivocab\_uncased}) & Retrieval semantico \\
    Indice vettoriale      & FAISS \texttt{IndexFlatIP}                     & ANN search \\
    Retrieval lessicale    & BM25Okapi (\texttt{rank\_bm25})                & Keyword search \\
    Fusione                & Reciprocal Rank Fusion ($k=60$)                & Re-ranking ibrido \\
    Generazione            & Phi-3.5-mini-instruct (3.8B)                   & Produzione risposta \\
    \bottomrule
  \end{tabular}
\end{table}

% ============================================================
\section{Implementazione del Retrieval}
% ============================================================

\subsection{Retrieval semantico}

\begin{pylisting}[title={Costruzione indice FAISS e retrieval semantico}]
def build_faiss_index(documents):
    embeddings = embedding_model.encode(
        documents, show_progress_bar=True,
        batch_size=64, convert_to_numpy=True,
        normalize_embeddings=True          # norma L2 = 1 (inner product = coseno)
    )
    idx = faiss.IndexFlatIP(embeddings.shape[1])  # dimensione 768
    idx.add(embeddings.astype(np.float32))
    return idx

def retrieve_semantic(query, faiss_index, documents, metadata, k=10):
    qe = embedding_model.encode(
        [query], convert_to_numpy=True, normalize_embeddings=True
    ).astype(np.float32)
    sims, idxs = faiss_index.search(qe, k)
    return [
        {'rank': i+1, 'score': float(s), 'idx': int(ix),
         'document': documents[ix], 'metadata': metadata[ix]}
        for i, (s, ix) in enumerate(zip(sims[0], idxs[0]))
    ]
\end{pylisting}

Il retrieval semantico codifica la query con SciBERT e recupera i $k$ documenti con massima similarit\`a coseno. FAISS \texttt{IndexFlatIP} calcola l'inner product su tutti i vettori: essendo normalizzati, coincide esattamente con la similarit\`a coseno.

\subsection{Retrieval BM25}

\begin{pylisting}[title={Costruzione indice BM25 e retrieval lessicale}]
def build_bm25_index(documents):
    tokenized = [doc.lower().split() for doc in documents]
    return BM25Okapi(tokenized), tokenized

def retrieve_bm25(query, bm25, documents, metadata, k=10):
    scores = bm25.get_scores(query.lower().split())
    top    = np.argsort(scores)[::-1][:k]
    return [
        {'rank': i+1, 'score': float(scores[ix]), 'idx': int(ix),
         'document': documents[ix], 'metadata': metadata[ix]}
        for i, ix in enumerate(top)
    ]
\end{pylisting}

BM25 tokenizza i documenti con semplice split su spazio e punteggia ogni documento in base alla frequenza normalizzata dei termini della query. Non richiede inferenza neurale, quindi \`e ordini di grandezza pi\`u veloce del retrieval semantico a runtime.

% ============================================================
\section{Generazione della Risposta}
% ============================================================

\begin{pylisting}[title={Prompt e generazione con Phi-3.5-mini-instruct}]
def rag_answer(query, results):
    context_parts = [
        f"[Paper {r['rank']}] {r['metadata']['title']}\n{r['document']}"
        for r in results
    ]
    context = "\n\n".join(context_parts)
    if len(context) > 3000:
        context = context[:3000] + "\n[...truncated...]"

    messages = [
        {"role": "system", "content": (
            "You are a helpful assistant specialized in NLP research. "
            "Answer questions based only on the provided research papers. "
            "Always cite the paper numbers [Paper X] you used. "
            "If the context is insufficient, say so clearly."
        )},
        {"role": "user", "content": (
            f"Context from research papers:\n\n{context}\n\n"
            f"Question: {query}\n\n"
            f"Answer based on the context above, citing paper numbers:"
        )}
    ]
    prompt = tokenizer.apply_chat_template(
        messages, tokenize=False, add_generation_prompt=True
    )
    output    = llm(prompt)
    full_text = output[0]['generated_text']
    return full_text[len(prompt):].strip()
\end{pylisting}

Il prompt \`e strutturato con ruolo \texttt{system} (vincola l'LLM all'uso del solo contesto fornito e richiede citazioni esplicite) e ruolo \texttt{user} (context + query). Il contesto \`e troncato a 3\,000 caratteri per gestire la finestra di contesto del modello su CPU. La temperatura \`e $T=0.3$ per favorire risposte deterministiche.

% ============================================================
\section{Metriche di Valutazione: Implementazione}
% ============================================================

\begin{pylisting}[title={Context Relevance e Answer Relevance}]
def context_relevance(query, results):
    qe   = embedding_model.encode(
               [query], convert_to_numpy=True, normalize_embeddings=True)
    snip = [r['document'][:200] for r in results]
    de   = embedding_model.encode(
               snip, convert_to_numpy=True, normalize_embeddings=True)
    return float(np.mean(de @ qe.T))   # media dei coseni query-snippet

def answer_relevance(query, answer):
    embs = embedding_model.encode(
               [query, answer], convert_to_numpy=True, normalize_embeddings=True)
    return float(embs[0] @ embs[1])    # coseno query-risposta
\end{pylisting}

\begin{pylisting}[title={Faithfulness proxy e robustezza}]
def faithfulness_proxy(answer, results):
    context_text  = ' '.join(r['document'][:500] for r in results).lower()
    words         = re.findall(r'\b[a-z]{4,}\b', answer.lower())
    content_words = [w for w in words if w not in sw]  # sw = stopwords NLTK
    if not content_words:
        return 0.0
    return sum(1 for w in content_words if w in context_text) / len(content_words)

def robustness_test(query, paraphrases, faiss_index, bm25,
                    documents, metadata, k=5):
    base_titles = {r['metadata']['title']
                   for r in retrieve_hybrid(
                       query, faiss_index, bm25, documents, metadata, k)}
    overlaps = []
    for para in paraphrases:
        para_titles = {r['metadata']['title']
                       for r in retrieve_hybrid(
                           para, faiss_index, bm25, documents, metadata, k)}
        overlaps.append(len(base_titles & para_titles) / k)
    return float(np.mean(overlaps))
\end{pylisting}

La \textbf{faithfulness proxy} \`e una stima lessicale della fedelt\`a: controlla che le parole contenutistiche della risposta siano ancorate al testo dei documenti recuperati. Non misura la correttezza logica (richiederebbe un NLI model), ma individua risposte ``fuori contesto''.

La \textbf{robustezza} testa tre query con due parafrasature ciascuna, misurando l'overlap@5 tra i risultati della query originale e quelli delle varianti. Un overlap basso segnala instabilit\`a del retriever a piccole variazioni lessicali.

% ============================================================
\section{Esperimento: Variazione di N\_DOCS}
% ============================================================

\subsection{Risultati per singolo N\_DOCS}

Le tabelle seguenti riportano i risultati completi per ciascuna configurazione (dati reali, esecuzione corrente con subset annidati e metodologia corretta, si veda la nota in Sezione~\ref{sec:nota-metodologica}).

\begin{table}[H]
  \centering
  \caption{Risultati N\_DOCS = 1\,000}
  \begin{tabular}{lrrr}
    \toprule
    \textbf{Metrica} & \textbf{Base} & \textbf{Ibrido} & \textbf{Delta} \\
    \midrule
    Precision@5        & 0.160 & 0.480 & +0.320 \\
    Context Relevance  & 0.716 & 0.714 & $-$0.002 \\
    Answer Relevance   & 0.676 & 0.657 & $-$0.019 \\
    Faithfulness       & 0.568 & 0.765 & +0.197 \\
    Robustezza         & 0.233 & 0.267 & +0.033 \\
    \bottomrule
  \end{tabular}
\end{table}

\begin{table}[H]
  \centering
  \caption{Risultati N\_DOCS = 5\,000}
  \begin{tabular}{lrrr}
    \toprule
    \textbf{Metrica} & \textbf{Base} & \textbf{Ibrido} & \textbf{Delta} \\
    \midrule
    Precision@5        & 0.280 & 0.680 & +0.400 \\
    Context Relevance  & 0.724 & 0.722 & $-$0.002 \\
    Answer Relevance   & 0.676 & 0.659 & $-$0.017 \\
    Faithfulness       & 0.613 & 0.719 & +0.106 \\
    Robustezza         & 0.167 & 0.100 & $-$0.067 \\
    \bottomrule
  \end{tabular}
\end{table}

\begin{table}[H]
  \centering
  \caption{Risultati N\_DOCS = 15\,000}
  \begin{tabular}{lrrr}
    \toprule
    \textbf{Metrica} & \textbf{Base} & \textbf{Ibrido} & \textbf{Delta} \\
    \midrule
    Precision@5        & 0.200 & 0.560 & +0.360 \\
    Context Relevance  & 0.713 & 0.731 & +0.018 \\
    Answer Relevance   & 0.683 & 0.675 & $-$0.008 \\
    Faithfulness       & 0.655 & 0.755 & +0.100 \\
    Robustezza         & 0.100 & 0.067 & $-$0.033 \\
    \bottomrule
  \end{tabular}
\end{table}

\begin{table}[H]
  \centering
  \caption{Risultati N\_DOCS = 25\,000}
  \begin{tabular}{lrrr}
    \toprule
    \textbf{Metrica} & \textbf{Base} & \textbf{Ibrido} & \textbf{Delta} \\
    \midrule
    Precision@5        & 0.200 & 0.560 & +0.360 \\
    Context Relevance  & 0.720 & 0.728 & +0.008 \\
    Answer Relevance   & 0.692 & 0.667 & $-$0.025 \\
    Faithfulness       & 0.628 & 0.722 & +0.094 \\
    Robustezza         & 0.067 & 0.000 & $-$0.067 \\
    \bottomrule
  \end{tabular}
\end{table}

\begin{table}[H]
  \centering
  \caption{Risultati N\_DOCS = 44\,949 (corpus completo)}
  \begin{tabular}{lrrr}
    \toprule
    \textbf{Metrica} & \textbf{Base} & \textbf{Ibrido} & \textbf{Delta} \\
    \midrule
    Precision@5        & 0.040 & 0.480 & +0.440 \\
    Context Relevance  & 0.719 & 0.724 & +0.005 \\
    Answer Relevance   & 0.687 & 0.680 & $-$0.007 \\
    Faithfulness       & 0.496 & 0.669 & +0.173 \\
    Robustezza         & 0.033 & 0.033 & 0.000 \\
    \bottomrule
  \end{tabular}
\end{table}

\subsection{Riepilogo comparativo: Precision@5}

\begin{table}[H]
  \centering
  \caption{Precision@5 al variare di N\_DOCS}
  \begin{tabular}{rrrrr}
    \toprule
    \textbf{N\_DOCS} & \textbf{Base} & \textbf{Ibrido} & \textbf{Delta} & \textbf{Miglioramento} \\
    \midrule
    1\,000  & 0.160 & 0.480 & +0.320 & +200.0\% \\
    5\,000  & 0.280 & 0.680 & +0.400 & +142.9\% \\
    15\,000 & 0.200 & 0.560 & +0.360 & +180.0\% \\
    25\,000 & 0.200 & 0.560 & +0.360 & +180.0\% \\
    44\,949 & 0.040 & 0.480 & +0.440 & +1100.0\% \\
    \bottomrule
  \end{tabular}
\end{table}

\subsection{Dettaglio per query: Precision@5, base vs.\ ibrido}

\begin{table}[H]
  \centering
  \caption{P@5 per singola query al variare di N\_DOCS -- Base / Ibrido}
  \small
  \begin{tabular}{p{6cm}rrrrr}
    \toprule
    \textbf{Query} & \textbf{1K} & \textbf{5K} & \textbf{15K} & \textbf{25K} & \textbf{44.9K} \\
    \midrule
    BERT / self-attention  & 0.0/0.6 & 0.2/0.4 & 0.0/0.4 & 0.2/0.4 & 0.0/0.4 \\
    word2vec               & 0.2/0.4 & 0.2/0.8 & 0.4/0.6 & 0.4/0.4 & 0.0/0.4 \\
    machine translation    & 0.2/0.4 & 0.4/0.8 & 0.4/0.6 & 0.2/0.8 & 0.0/0.4 \\
    topic modeling         & 0.0/0.4 & 0.0/0.4 & 0.0/0.4 & 0.0/0.4 & 0.0/0.4 \\
    named entity recog.    & 0.4/0.6 & 0.6/1.0 & 0.2/0.8 & 0.2/0.8 & 0.2/0.8 \\
    \midrule
    \textbf{Media}         & \textbf{0.16/0.48} & \textbf{0.28/0.68} & \textbf{0.20/0.56} & \textbf{0.20/0.56} & \textbf{0.04/0.48} \\
    \bottomrule
  \end{tabular}
\end{table}

\subsection{Robustezza al variare di N\_DOCS}

\begin{table}[H]
  \centering
  \caption{Robustezza (overlap@5 su query parafrasate): base vs.\ ibrido}
  \begin{tabular}{rrr}
    \toprule
    \textbf{N\_DOCS} & \textbf{Base} & \textbf{Ibrido} \\
    \midrule
    1\,000  & 0.233 & 0.267 \\
    5\,000  & 0.167 & 0.100 \\
    15\,000 & 0.100 & 0.067 \\
    25\,000 & 0.067 & 0.000 \\
    44\,949 & 0.033 & 0.033 \\
    \bottomrule
  \end{tabular}
\end{table}

% ============================================================
\section{Risultati e Discussione}
% ============================================================

\subsection{Il retrieval ibrido migliora sistematicamente la Precision@5}

Il dato pi\`u netto \`e che il sistema ibrido supera il sistema base in \textit{tutte e cinque} le configurazioni di N\_DOCS, con miglioramenti che vanno da +142.9\% a +1100\%. L'ibrido porta la Precision@5 media da valori tra 0.040 e 0.280 a valori tra 0.480 e 0.680: il salto \`e ancora pi\`u netto rispetto a quanto misurato nella prima esecuzione (Sezione~\ref{sec:nota-metodologica}), perch\'e ora il confronto \`e fatto a parit\`a di corpus fra le due configurazioni.

Questo conferma l'ipotesi di partenza: BM25, catturando i termini tecnici esatti (es.\ ``NER'', ``LDA'', ``skip-gram''), recupera documenti che il solo retrieval semantico manca, perch\'e SciBERT tende a generalizzare i termini tecnici verso concetti pi\`u astratti. La fusione RRF bilancia i punti di forza dei due retriever.

\subsection{Il crollo del sistema base sul corpus completo}

Il risultato pi\`u netto di questa esecuzione \`e il crollo della Precision@5 base a \textbf{0.040} sull'intero corpus (44\,949 documenti): guardando il dettaglio per query (Tabella~7), a quella dimensione il retrieval puramente semantico non trova \emph{nessun} documento rilevante nei top-5 per 4 query su 5 (BERT/self-attention, word2vec, machine translation, topic modeling azzerate), e recupera qualcosa solo per ``named entity recognition'' (0.2). Il sistema ibrido, sullo stesso corpus, resta invece a \textbf{0.480} -- pressoch\'e invariato rispetto alle configurazioni pi\`u piccole. Questo \`e ora un confronto pulito: gli stessi 44\,949 documenti, la stessa query, lo stesso gold standard, l'unica differenza \`e la presenza di BM25 nella fusione. Il retrieval solo-embedding, su un corpus molto ampio e con query brevi e generiche, perde quasi completamente la capacit\`a di isolare i pochi documenti realmente pertinenti tra le decine di migliaia disponibili; BM25 mantiene invece il segnale sulle parole chiave indipendentemente da quanti documenti competono nell'indice.

\subsection{Andamento non monotono della Precision@5 con N\_DOCS}

Sia la curva base sia quella ibrida hanno un picco a N=5\,000 (0.280 e 0.680) e calano dopo: il corpus a 5\,000 documenti \`e abbastanza grande da contenere paper rilevanti, ma abbastanza piccolo da non essere ``diluito'' da troppi documenti distanti dal tema della query. Il calo successivo \`e per\`o molto pi\`u marcato per il sistema base (fino a 0.040) che per l'ibrido (che si stabilizza a 0.480--0.560): la fusione con BM25 non elimina l'effetto del corpus grande, ma ne assorbe gran parte.

\subsection{Context Relevance e Answer Relevance: differenze minime}

La Context Relevance (CR) e l'Answer Relevance (AR) mostrano differenze piccole tra base e ibrido (tipicamente $|\Delta| < 0.02$, tranne un paio di casi fino a $0.025$). Questo \`e atteso: entrambe le metriche sono calcolate con gli stessi embeddings SciBERT usati anche per il retrieval semantico, quindi condividono in parte lo stesso ``punto di vista'' -- non sono una misura indipendente della qualit\`a del sistema, un limite gi\`a dichiarato in partenza e confermato dai dati.

L'AR ibrida \`e quasi sempre leggermente inferiore a quella base: i documenti trovati da BM25 sono meno ``vicini'' alla query nello spazio degli embedding (li trova per parole chiave, non per significato), quindi il contesto che l'LLM riceve, pur pi\`u corretto nel contenuto, pu\`o produrre risposte con embedding leggermente pi\`u lontani dalla query.

\subsection{Faithfulness: il sistema ibrido produce risposte pi\`u fedeli}

La Faithfulness \`e sistematicamente pi\`u alta per il sistema ibrido in tutte le configurazioni, con un margine molto pi\`u ampio di quanto misurato in precedenza ($+0.094$ a $+0.197$, contro $+0.027$/$+0.095$ della prima esecuzione). Il salto pi\`u grande \`e proprio a N=1\,000 ($+0.197$) e sul corpus completo ($+0.173$): i documenti trovati dall'ibrido restano pi\`u legati al vocabolario tecnico della domanda, quindi l'LLM usa pi\`u parole che compaiono davvero nei testi forniti, invece di ``riempire i vuoti'' con parole plausibili ma non presenti nel contesto.

\subsection{Robustezza: non \`e pi\`u una vittoria automatica dell'ibrido}

Con il dato del sistema base ora disponibile per il confronto (Tabella~8), la robustezza racconta una storia pi\`u sfumata di quanto sembrasse: l'ibrido \`e pi\`u robusto del base solo a N=1\,000 (0.267 contro 0.233); a N=5\,000, 15\,000 e 25\,000 \`e il sistema \textbf{base} a risultare pi\`u stabile alle parafrasi (es.\ 0.167 contro 0.100 a N=5\,000), e a N=44\,949 i due sistemi sono pari (0.033). In altre parole, aggiungere BM25 non stabilizza automaticamente il retrieval rispetto a piccole variazioni della formulazione della query -- anzi, introducendo un secondo canale di recupero basato su corrispondenza esatta delle parole, l'ibrido pu\`o diventare pi\`u sensibile a un sinonimo diverso nella parafrasi, che BM25 semplicemente non riconosce come la stessa parola chiave. In entrambi i sistemi, comunque, la robustezza cala con corpus pi\`u grandi: con pi\`u documenti candidati, piccole variazioni lessicali della query spostano pi\`u facilmente quali documenti finiscono nei top-5.

\subsection{Query con performance diverse}

La query ``named entity recognition'' resta quella con P@5 pi\`u alta e stabile in versione ibrida (0.6--1.0): il termine ``NER'' \`e molto specifico e BM25 lo gestisce perfettamente. La query ``topic modeling'' \`e invece la pi\`u debole in assoluto: P@5 base sempre a 0.0 e ibrida ferma a 0.4 in ogni configurazione, segno che n\'e il retrieval semantico n\'e BM25 trovano bene i pochi paper specificamente su LDA/BERTopic in questo corpus. La query su BERT resta un caso instabile anche nella versione ibrida (0.4 quasi ovunque, mai sopra 0.6): un limite osservato concretamente anche nelle risposte generate (si veda l'esempio sotto).

\begin{table}[H]
\centering
\begin{tabular}{p{14.5cm}}
\toprule
\textbf{Esempio -- risposta generata sulla query BERT, corpus completo (44\,949 doc.)} \\
\midrule
\textbf{Base}: ``Based on the provided research papers, there is no direct information or discussion about BERT (Bidirectional Encoder Representations from Transformers) or self-attention mechanisms. The papers focus on various aspects of Natural Language Processing (NLP)\ldots'' \\[0.5em]
\textbf{Ibrido}: ``Based on the provided research papers, there is no direct explanation of BERT or self-attention mechanisms. The papers discuss various topics related to natural language processing (NLP), such as attention in neural machine translation (Paper 2), cross-lingual syntactic ability of multilingual BERT (Paper 4)\ldots'' \\
\bottomrule
\end{tabular}
\end{table}

In entrambi i casi l'LLM dichiara correttamente l'assenza di un paper specificamente dedicato a BERT invece di inventare una risposta -- un comportamento di fedelt\`a al contesto corretto, ma che conferma il limite: nemmeno l'ibrido, su questo corpus e con questo gold standard, recupera un paper che spieghi BERT direttamente; l'ibrido arriva per\`o pi\`u vicino al tema (cita esplicitamente un paper su BERT multilingue).

\subsection{Nota metodologica: correzione applicata}
\label{sec:nota-metodologica}

Una prima esecuzione di questo esperimento presentava alcuni limiti metodologici, poi corretti nel codice (\texttt{lab5rag.py} e \texttt{lab5\_hybrid.py}) prima di generare i numeri riportati sopra:

\begin{itemize}
  \item \textbf{Campionamento non annidato}: ogni corpus a $N$\_DOCS diverso era ottenuto con una chiamata indipendente a \texttt{df.sample(...)}, quindi un paper poteva comparire nel corpus da 5\,000 documenti e non in quello da 44\,949 (o viceversa) per puro caso, non per un vero effetto della dimensione del corpus. Corretto mischiando il dataset una sola volta e facendo prendere a ogni corpus i primi $N$ documenti di quell'ordine fisso, cos\`i che i corpus pi\`u piccoli siano sempre sottoinsiemi di quelli pi\`u grandi.
  \item \textbf{Precision@k non uniforme fra i due script}: uno cercava le parole chiave in titolo+abstract, l'altro solo nel titolo. Corretto usando lo stesso criterio in entrambi.
  \item \textbf{Corrispondenza troppo permissiva}: parole chiave come ``ner'' o ``bert'' potevano matchare per errore parole diverse (``corner'', ``RoBERTa''). Corretto cercando la parola (o frase) intera.
  \item \textbf{Generazione non deterministica}: veniva usato il campionamento casuale nella generazione, quindi risposte e metriche potevano cambiare a ogni esecuzione. Corretto con generazione greedy (sempre la parola pi\`u probabile), riproducibile.
\end{itemize}

Il fix pi\`u importante \`e il primo: senza subset annidati, la curva ``Precision@5 al variare di N\_DOCS'' non era interpretabile, perch\'e un calo poteva dipendere solo da quali paper erano finiti per caso in un campione. Con i subset annidati, il crollo del sistema base a N=44\,949 (Sezione~11.2) \`e ora un risultato solido: lo stesso identico corpus, con e senza BM25, produce una differenza di quasi 12 volte in Precision@5.

% ============================================================
\section{Conclusioni}
% ============================================================

Il laboratorio ha implementato un sistema RAG ibrido su un corpus di paper scientifici NLP, dimostrando che la combinazione di retrieval semantico (SciBERT + FAISS) e lessicale (BM25) tramite Reciprocal Rank Fusion porta a miglioramenti consistenti e ora misurati in modo pulito (subset annidati, criteri di valutazione uniformi, generazione deterministica).

I risultati principali sono:

\begin{enumerate}
  \item \textbf{Il retrieval ibrido migliora la Precision@5 in modo sistematico}: il delta medio sulle cinque configurazioni \`e +0.376, con un miglioramento percentuale che va da +142.9\% a +1100\%.

  \item \textbf{Il sistema base crolla su corpus grandi, l'ibrido no}: a 44\,949 documenti la Precision@5 base scende a 0.040 (quasi zero), mentre l'ibrido resta a 0.480 -- la differenza pi\`u importante misurata in tutto l'esperimento, e la ragione pratica per cui BM25 non \`e opzionale su corpus di queste dimensioni.

  \item \textbf{La Faithfulness beneficia chiaramente del retrieval ibrido}: margini fra $+0.094$ e $+0.197$ in tutte le configurazioni, il risultato pi\`u stabile di tutto l'esperimento.

  \item \textbf{La robustezza alle parafrasi non premia sempre l'ibrido}: solo a N=1\,000 l'ibrido \`e pi\`u stabile del base; a N=5\,000--25\,000 \`e vero il contrario. Aggiungere BM25 aiuta la precisione ma non garantisce automaticamente pi\`u stabilit\`a a piccole variazioni della query.
\end{enumerate}

I limiti principali riguardano il gold standard (definito manualmente e limitato a 5 query), il costo computazionale su CPU (necessario in assenza di GPU, e rilevante: l'intero esperimento su 5 configurazioni richiede diverse ore di generazione), e la faithfulness proxy che \`e una stima lessicale e non valuta la correttezza logica delle risposte.

\end{document}
